\documentclass[11pt]{article}
\usepackage{amsmath,amssymb,amsthm}
\newtheorem{lemma}{Lemma}
\newtheorem{theorem}{Theorem}
\title{Positive common root pools supplied by auxiliary routing}
\author{}\date{}
\begin{document}\maketitle

\paragraph{Status.}
This is a complete written root-only compatibility theorem, with exact
rational checks. It permits a single strictly positive residue pool for
several prescribed words by choosing routing at unused query indices.
It is not Lean-checked, does not preserve descendant inequalities, and
does not prove a growing-query compiler.

Let $q\geq1$, and let $t_1<\cdots<t_q$ be exterior queries. Write
$v(p)=(1/(t_j-p))_{j=1}^q$ for $p<t_1$.

\begin{lemma}[Positive spanning from alternating Cauchy columns]
For any $q+1$ distinct poles $p_0<\cdots<p_q<t_1$, the vectors
$(-1)^h v(p_h)$ positively span $\mathbb R^q$. More precisely,
every $y\in\mathbb R^q$ is a linear combination of these vectors
with all coefficients strictly positive.
The same conclusion holds after adding any other signed columns
and prescribing arbitrary strictly positive coefficients on them.
\end{lemma}
\begin{proof}
Every $q$ by $q$ minor of the Cauchy matrix $(1/(t_j-p_h))$ is
nonzero and has the same sign, by the Cauchy determinant formula:
its factors are the ordered differences of queries and poles,
with positive denominator factors. The standard signed-minor
kernel vector consequently has alternating nonzero coordinates.
After the column sign changes $(-1)^h$, this gives a strictly
positive kernel vector $z$. The signed matrix has rank $q$.
For an arbitrary $y$, take any real solution $c$ of $Ac=y$.
Then $c+\lambda z$ is strictly positive for sufficiently large
$\lambda$, without changing its image. Additional prescribed
columns can first be subtracted from $y$.
\end{proof}

Suppose now that the poles come in $M$ adjacent pairs,
\[
 p_1^-<p_1^+<p_2^-<p_2^+<\cdots<p_M^-<p_M^+<t_1.
                                                               \tag{1}
\]
A routing word $\theta\in\{-1,1\}^M$ and positive pool
$c=(c_s^-,c_s^+)_{s=1}^M$ give the root score
\[
 F_{\theta,c}(t)=\sum_{s=1}^M\theta_s
       \left(\frac{c_s^-}{t-p_s^-}
                     -\frac{c_s^+}{t-p_s^+}\right).            \tag{2}
\]
Let $I\subseteq\{1,\ldots,M\}$ have size $q$. Indices in $I$ are
prescribed routing indices; the others are available for auxiliary
routing. An association between queries and $I$ may be fixed in
any order.

\begin{theorem}[Several words with one positive pool]
Let $K\geq1$ and set $L=\lceil(q+1)/2\rceil$. Assume
\[
                         M-q\geq KL.                          \tag{3}
\]
For every collection of prescribed routing words
$\rho^{(a)}\in\{-1,1\}^I$ and every collection of desired real
query vectors $y^{(a)}\in\mathbb R^q$, $1\leq a\leq K$, there
exist completions $\theta^{(a)}\in\{-1,1\}^M$ and one common
strictly positive pool $c$ such that
\[
 \theta^{(a)}|_I=\rho^{(a)},\qquad
 F_{\theta^{(a)},c}(t_j)=y_j^{(a)}
                   \quad(1\leq a\leq K,\ 1\leq j\leq q).
                                                               \tag{4}
\]
In particular arbitrary prescribed root sign words can be realized
by setting $y_j^{(a)}$ equal to their requested signs.
\end{theorem}
\begin{proof}
Choose $K$ linearly independent sign vectors
$w_1,\ldots,w_K\in\{-1,1\}^K$. For example take $w_1$ to be
all positive, and for $2\leq h\leq K$ obtain $w_h$ by flipping
coordinate $h-1$ of $w_1$. Subtracting the first column from the
others proves independence, including the case $K=1$.

Choose disjoint groups $J_1,\ldots,J_K$ of $L$ auxiliary pair
indices each. For every $s\in J_h$, set
$\theta_s^{(a)}=(w_h)_a$ for all $a$. Set any remaining auxiliary
routing bits arbitrarily. Prescribe positive residues, for example
one, on the prescribed pairs and on those remaining auxiliary
pairs. Subtract their score vectors from the desired matrix
$Y=(y_j^{(a)})$. Since the columns $w_h$ form a basis, the remaining
matrix has a unique decomposition
\[
                           \sum_{h=1}^K w_h z_h^{\mathsf T}
                \quad\text{with }z_h\in\mathbb R^q.
\]
Within group $J_h$, sorted pole columns alternate in sign in the
scalar transform in parentheses in (2). There are $2L\geq q+1$
such columns. The lemma gives strictly positive coefficients on
all of them whose scalar query vector is exactly $z_h$.
Their contributions to the full $K$ by $q$ matrix are therefore
$w_h z_h^{\mathsf T}$. Combining all groups proves (4).
\end{proof}

For two target words condition (3) permits approximately half of the
pair indices to be queried. This conclusion differs from a common-pool
assertion with all routing bits fixed: the auxiliary completions are
chosen separately for the different target words. In the dyadic pole
family $p_s^-=p_s$, $p_s^+=p_s+\delta$, with $\delta$ smaller than
every gap between distinct $p_s$, the adjacency hypothesis (1) holds.
The theorem thus applies to that family's root score whenever its
unused routing bits can be chosen independently.

\paragraph{Why this does not yet recurse.}
The common pool in the theorem need not satisfy a single additional
descendant query. The proof uses exact value interpolation, and its
combined query matrix has rank $Kq$; it obtains this rank by reserving
about $Kq$ pole coefficients. Stacking all descendant rows would
exceed that supply once the desired total number of queries grows
faster than the number of poles. Also, an induction requiring common
pools for $K$ full subtrees generally doubles the number of subtrees
when passing to their children. Condition (3) by itself then yields
only a bounded total query-to-pole ratio.

An actual superconstant improvement would require a weaker invariant,
such as intersections of projected sign cones, that retains these
auxiliary choices without matching all descendant values. Neither the
positive-spanning lemma nor the count of possible completions proves
that stronger statement.

\paragraph{Exact verification.}
\texttt{auxiliary-root-pool-check.py} uses rational adjacent poles and
exterior queries, computes signed Cauchy kernels and inverses exactly,
and constructs positive pools by a kernel shift. It checks every
collection of root sign words for $1\leq q,K\leq3$, and for
$(q,K)=(4,2),(5,2)$, at the boundary value $M=q+KL$.
All original query scores equal their desired signs exactly.
No floating-point feasibility test or descendant assertion is used.
\end{document}
